% =========================================================
\chapter{Simulações, Análises e Discussão de Resultados}
% =========================================================

Este capítulo deve apresentar evidências capazes de sustentar as afirmações
científicas do trabalho. Resultado não deve ser tratado simplesmente como uma
figura, tabela ou valor numérico: cada resultado precisa ser interpretado.

% ---------------------------------------------------------
\section{Considerações iniciais}
% ---------------------------------------------------------

\begin{enumerate}
	\item Qual é a finalidade dos experimentos ou simulações?
	\item Quais objetivos específicos estão sendo avaliados?
	\item Quais hipóteses estão sendo testadas?
	\item Quais métricas serão utilizadas?
	\item Quais resultados são esperados?
	\item Como os resultados serão comparados?
	\item Qual é o critério de sucesso?
	\item Como a validade dos resultados será assegurada?
	\item Quais fontes de incerteza existem?
	\item Como os resultados serão discutidos?
\end{enumerate}

% ---------------------------------------------------------
\section{Cenários, métodos e parâmetros de simulação/teste/coleta de dados}
% ---------------------------------------------------------

\begin{enumerate}
	\item Qual ambiente foi utilizado?
	\item Qual software, hardware ou plataforma foi utilizada?
	\item Qual versão das ferramentas foi empregada?
	\item Quais parâmetros foram utilizados?
	\item Quais valores foram mantidos constantes?
	\item Quais parâmetros foram variados?
	\item Quais cenários foram considerados?
	\item Quantas repetições foram realizadas?
	\item Como os dados foram coletados e tratados?
	\item Como outro pesquisador poderia reproduzir o experimento?
\end{enumerate}

{Pergunta fundamental.}
\textbf{Um pesquisador externo conseguiria reproduzir exatamente este
	experimento apenas lendo esta seção?}

Se a resposta for negativa, a metodologia ainda não está suficientemente
documentada.

% ---------------------------------------------------------
\section{Apresentação e discussão dos resultados}
% ---------------------------------------------------------

Para cada resultado relevante, o estudante deve responder às seguintes
perguntas:

\begin{enumerate}
	\item O que está sendo mostrado?
	\item Por que esse resultado é importante?
	\item Qual variável está sendo analisada?
	\item Qual tendência pode ser observada?
	\item Por que essa tendência ocorre?
	\item O resultado confirma ou contradiz a hipótese?
	\item O resultado está de acordo com a literatura?
	\item Por que o resultado difere ou coincide com trabalhos anteriores?
	\item Qual é a magnitude da melhoria ou degradação observada?
	\item Qual conclusão científica pode ser extraída desse resultado?
\end{enumerate}

{Estrutura recomendada para discussão.}

\begin{center}
	\textbf{
		Resultado
		$\rightarrow$
		comportamento observado
		$\rightarrow$
		explicação física/matemática
		$\rightarrow$
		comparação
		$\rightarrow$
		interpretação
		$\rightarrow$
		conclusão.
	}
\end{center}

% ---------------------------------------------------------
\section{Considerações finais}
% ---------------------------------------------------------

\begin{enumerate}
	\item Quais foram os principais resultados?
	\item Quais objetivos específicos foram atendidos?
	\item Quais hipóteses foram confirmadas?
	\item Quais hipóteses foram rejeitadas?
	\item Qual foi o melhor cenário?
	\item Quais parâmetros mais influenciaram o desempenho?
	\item Quais limitações foram identificadas?
	\item Os resultados são estatisticamente ou tecnicamente significativos?
	\item Os resultados são coerentes com a literatura?
	\item Quais evidências permitem afirmar que a solução funciona?
\end{enumerate}
